\chapter{Glossary}\label{app:glossary}

This glossary explains the technical words used in this manual. Each word is
also printed in \textbf{\textit{bold italics}} where the manual first uses it.
Where a word has more than one meaning, the meanings are numbered.

\begingroup
\setlist[description]{font=\normalfont\bfseries,leftmargin=1.5em,itemsep=3pt,parsep=0pt}

\section*{0--9}

\begin{description}
\item[12-6-4 ion model] A more accurate set of parameters for metal ions that
  adds an extra attraction term. ProPrep adds it to the topology with ParmEd.
  See Chapter~\ref{ch:topology}.
\item[2mFo--DFc map] The electron-density map into which an X-ray model was
  built. See Chapter~\ref{ch:compare}.
\end{description}

\section*{A}

\begin{description}
\item[accession] The code that identifies a protein in UniProt, for example
  P69905 (human hemoglobin alpha).
\item[ACE] An acetyl cap: a small neutral group added before the first residue
  of a chain, so the chain start carries no artificial charge.
\item[alpha carbons (CA)] The central carbon atom of each amino acid. They give
  one point per residue and are often used to compare protein shapes.
\item[AlphaFill] A database that takes AlphaFold models and copies in
  (``transplants'') ligands, cofactors and metal ions from similar
  experimental structures.
\item[AlphaFold] A program that predicts a protein's 3D structure from its
  sequence. Its models are available from the AlphaFold Database by UniProt
  accession. See Chapter~\ref{ch:structure}.
\item[alternate locations (altlocs)] Atoms that the experiment saw in two or
  more positions, each for a share of the time (its occupancy). A simulation
  needs exactly one position per atom.
\item[AM1-BCC] A fast way to compute partial charges for a small molecule,
  built into AmberTools. Less careful than RESP, but it needs no
  quantum-chemistry program. See Chapter~\ref{ch:forcefields}.
\item[AMBERHOME] A setting in the Terminal that names the folder of an Amber
  installation. If it points to another Amber, ProPrep uses that one instead of
  its own. See Chapter~\ref{ch:installation}.
\item[AmberTools] The free part of the Amber simulation package (tLEaP,
  sander, cpptraj, antechamber, ParmEd and others). It is installed together
  with ProPrep.
\item[AS4] The Amber name for an aspartate that may gain or lose a proton
  during a constant-pH simulation.
\item[asymmetric unit] The chains actually written in a crystal-structure
  file. It may hold part of the working molecule, exactly one copy, or
  several copies.
\item[atom type] A short code that tells the force field what kind of atom
  this is, for example \texttt{CX} for a protein C-alpha. Bond, angle and other
  parameters are looked up by atom type.
\item[axial ligand] A protein side chain that binds the metal of a flat ring,
  such as a heme, from above or below the ring. In cytochrome~c these are His18
  and Met80.
\end{description}

\section*{B}

\begin{description}
\item[backbone] The atoms N, CA, C and O that every amino acid has. Comparing
  backbones works even when the side chains differ.
\item[batch processing] Running a template over a list of proteins without
  typing. See Chapter~\ref{ch:batch}.
\item[biological assembly] The form of the molecule believed to exist and work
  in the cell, for example the hemoglobin tetramer. It is built from the
  asymmetric unit using the instructions in REMARK~350.
\item[BLAST] A program that compares a sequence with all sequences in a
  database and lists the most similar ones. ProPrep uses the public BLAST
  server of the NCBI.
\item[box padding] In a membrane system, the width of lipid kept between the
  protein and the side of the box. The default is 15\,Å.
\item[breadcrumb] The \texttt{Navigation:} line at the top of each ProPrep
  screen, which shows where you are in the menus.
\item[broken-symmetry] For iron--sulfur clusters: one of several equivalent
  ways of pairing the spins of the iron atoms. ProPrep's parameter sets come
  in several such variants (codes such as FO1, FR2, FD0).
\item[b-type heme] A heme held in the protein only by its axial ligands,
  without covalent bonds to the protein (PDB name HEM, as in hemoglobin).
\item[buffer] The smallest distance between the protein and the edge of the
  simulation box, usually 10\,Å.
\item[buried surface area (BSA)] The surface two chains hide from water by
  touching each other, in Å\textsuperscript{2}. Large values (over about
  1000\,Å\textsuperscript{2}) suggest a real complex.
\end{description}

\section*{C}

\begin{description}
\item[campaign] Docking every molecule of a library file into one receptor and
  box with the same settings, giving a results table and a best pose per
  molecule.
\item[cap (terminal cap)] A small neutral group (ACE or NME) that closes a
  chain end instead of the charged terminus. It is used when the real protein
  continues beyond the structure.
\item[chain break] A place where the protein chain continues in reality but
  the atoms are missing from the file. Such an end is left neutral; a real
  chain end is charged.
\item[charge (of a QM region)] The total charge of the atoms treated with
  quantum mechanics, always a whole number.
\item[chimera] An engineered protein made from parts of two proteins, for
  example the $\beta_2$-adrenergic receptor with T4 lysozyme
  inserted (2RH1).
\item[cluster profile] A saved description of one computing cluster (job
  system, modules, partitions, engines) that ProPrep uses to write job
  scripts.
\item[component] In the PDB Filter, one kind of material in a chain: the
  protein, the waters, or a particular ligand or ion.
\item[composition] What a PDB entry contains, in the depositors' words:
  molecule names, chains, organism and mutations.
\item[conda] A program that installs scientific software and everything it
  needs into separate, self-contained folders.
\item[conda environment] One such self-contained folder of software. The conda
  installer puts ProPrep in an environment named \texttt{ProPrep}, switched on
  with \texttt{conda activate ProPrep}.
\item[constant-pH simulation] A simulation in which chosen residues can gain
  and lose protons while it runs, so their states follow the pH instead of
  being fixed.
\item[coordinate file (coordinates)] The file (\file{.rst7}) that gives the
  position of every atom, and the box size. It belongs to one topology; after
  a simulation stage it also holds the atom velocities.
\item[counter-ions] Ions of opposite charge added to make the total charge of
  the simulation box zero.
\item[cpin] The Amber input file that lists the titratable residues and their
  possible states for a constant-pH simulation.
\item[C-terminus] The end of a protein chain, with a carboxyl group that
  normally carries a negative charge.
\item[c-type heme] A heme bonded to the protein through two cysteines (PDB
  name HEC, as in cytochrome~c).
\end{description}

\section*{D}

\begin{description}
\item[default] The answer shown in parentheses at the end of a question, for
  example \texttt{(y)}. Pressing \key{Enter} without typing anything accepts
  it.
\item[desolvation] The effect of a residue being buried away from water. It
  favors the uncharged form and so shifts the pKa.
\item[dihedral] A twist angle around a bond, defined by four atoms in a row.
\item[DISANG] An Amber file of restraints on distances, angles or twists
  between specific atoms.
\item[disulfide bond] A covalent bond between the sulfur atoms of two
  cysteines. In Amber, both cysteines must be named CYX and joined by a
  \texttt{bond} line. See Chapter~\ref{ch:topology}.
\item[docking] Predicting where a small molecule sits in a protein's pocket,
  and how tightly, by trying many positions and scoring each one. See
  Chapter~\ref{ch:docking}.
\item[DOI] Digital Object Identifier: a permanent web address for a
  publication, software release or data set.
\end{description}

\section*{E}

\begin{description}
\item[electron density] The experimental map of an X-ray structure into which
  the atomic model is built.
\item[electronic embedding] In QM/MM, the MM partial charges act on the QM
  electrons, so the protein polarizes the QM region.
\item[E-value] In BLAST, the number of hits this good you would expect by
  chance. The smaller it is, the more certain the relationship.
\item[exhaustiveness] How hard a docking search works. Higher is slower and
  more thorough; the default is 8.
\item[extension] The end of a file name after the last dot, such as
  \file{.pdb}. It tells you what kind of file it is. See
  Appendix~\ref{app:filetypes}.
\end{description}

\section*{F}

\begin{description}
\item[force field] The set of rules and numbers (parameters) that describes
  how atoms attract, repel and hold on to each other in a simulation, for
  example ff19SB for proteins. It recognizes residues by their names. See
  Chapter~\ref{ch:forcefields}.
\item[formal charge] The whole-number charge given to an atom or ion by
  chemistry rules (Zn\textsuperscript{2+}: +2, a thiolate sulfur: $-1$), as
  opposed to a partial charge.
\item[frame] One snapshot in a trajectory.
\item[Full Menu] ProPrep's main menu that shows every tool at once, each with
  a fixed code such as \key{1a} or \key{4e}. This manual uses Full Menu codes.
  See Chapter~\ref{ch:menus}.
\end{description}

\section*{G}

\begin{description}
\item[GAFF2] The General Amber Force Field, version~2: Amber's force field for
  drug-like organic molecules.
\item[Generalized Born (GB)] The most common implicit-solvent model in Amber.
\item[GL4] The Amber name for a glutamate that may gain or lose a proton
  during a constant-pH simulation.
\item[GROUP restraints] An older Amber format for restraining several groups of
  atoms, each with its own strength.
\end{description}

\section*{H}

\begin{description}
\item[HID] The Amber name for neutral histidine with its proton on the delta
  nitrogen (ND1).
\item[HIE] The Amber name for neutral histidine with its proton on the epsilon
  nitrogen (NE2); the more common neutral form.
\item[HIP] The Amber name for positively charged histidine, with protons on
  both ring nitrogens.
\item[homology model] A model of a protein built from the known structure of a
  related protein. ProPrep needs MODELLER to build one.
\end{description}

\section*{I}

\begin{description}
\item[identity] In a sequence alignment, the percentage of aligned positions
  that have the same amino acid.
\item[implicit solvent] Water replaced by a smooth surrounding medium instead
  of water molecules. Cheaper, but less realistic, than explicit water.
\item[isoelectric point (pI)] The pH at which a protein's net charge is zero.
\item[issue] A report (problem, question or suggestion) posted on a software
  project's public GitHub page. ProPrep's feedback form can create one.
\end{description}

\section*{L}

\begin{description}
\item[large model] In the MCPB metal-site method, the bigger cut-out of the
  metal site, used for the quantum calculation that gives the charges.
\item[layers (ONIOM)] The parts into which ONIOM splits the system: HIGH
  (quantum mechanics), an optional MEDIUM (a cheaper quantum method) and LOW
  (force field).
\item[leaflet] One of the two layers of a lipid bilayer. An asymmetric
  membrane has different lipids in each leaflet.
\item[license key] A personal code you receive after registering for a
  program, here MODELLER, that lets it run. See
  Chapter~\ref{ch:installation}.
\item[ligand] (1)~In docking, the small molecule being docked. (2)~For metals
  and cofactors, a protein residue (or water) whose atoms bind the metal and
  hold it in place.
\item[link atom] A hydrogen atom added where a bond crosses from the QM region
  to the MM region, so the QM region has no loose bond.
\item[lipid bilayer] The two-layer sheet of fat-like molecules that forms a
  cell membrane. Simulations of membrane proteins build a small patch of one.
\item[lipids] The fat-like molecules a membrane is made of. Each has a
  water-liking head and oily tails.
\item[localhost (127.0.0.1)] The name and address a computer uses for itself.
  A web page at this address comes from a program on your own computer, not
  from the internet.
\item[loops] (1)~Flexible parts of a protein chain between its helices and
  strands, often missing from experimental structures. (2)~In membrane
  building, rounds of rearranging molecules during packing.
\end{description}

\section*{M}

\begin{description}
\item[mask] A short Amber code that selects atoms, for example \texttt{:1-76}
  (residues 1 to 76) or \texttt{@CA} (all atoms named CA).
\item[MCPB] Metal Center Parameter Builder: the Amber method for making
  parameters for a metal and the residues that bind it, from quantum
  calculations on a small and a large model.
\item[MD engine] The program that actually runs the simulation:
  \texttt{sander} and \texttt{sander.MPI} (included with ProPrep), or
  \texttt{pmemd}, \texttt{pmemd.MPI} and \texttt{pmemd.cuda} (a separate Amber
  package). See Chapter~\ref{ch:runningmd}.
\item[mdin] The Amber input file that holds the settings of one simulation
  stage: number of steps, temperature, restraints and so on.
\item[mdout] Amber's text output for a stage: energies, temperature and
  pressure printed at regular intervals.
\item[mechanical embedding] In QM/MM, the MM charges do not act on the QM
  electrons; the two regions interact only through the force field.
\item[MEMEMBED] A program that works out how a protein sits in a membrane.
  ProPrep's default for orienting membrane proteins.
\item[mFo--DFc map] The difference electron-density map. It shows where the
  experimental data and the model disagree.
\item[microstate] One version of a protein that differs from the others only
  in the state of some sites, for example which hemes are oxidized or which
  disulfides are formed. See Chapter~\ref{ch:redox}.
\item[missing atoms] Atoms of a residue, usually side-chain atoms, that the
  experiment did not see.
\item[missing residues] Whole residues that were in the sample but not seen in
  the experiment, typically flexible loops and chain ends.
\item[MM (molecular mechanics)] Treating atoms as balls and springs with a
  force field, as in MD. Fast, but it cannot describe electrons.
\item[mmCIF] The newer file format for structures (\file{.cif} files), used by
  AlphaFold, AlphaFill and very large PDB entries. Most of ProPrep expects PDB
  format.
\item[MODELLER] A program from the Sali lab that builds missing parts of
  protein structures, standard mutations and homology models. It is installed
  with ProPrep but needs a free academic license key.
\item[molar ratio] How many molecules of one lipid are used for each of the
  others, for example POPC to cholesterol at 3:1.
\item[molecular dynamics (MD)] A calculation that moves every atom a tiny step
  at a time according to the forces between atoms, producing a film of how a
  molecule moves.
\item[MPI] A system for splitting one calculation between several processor
  cores, used by \texttt{sander.MPI} and \texttt{pmemd.MPI}.
\item[multiplicity] The number of unpaired electrons plus one: 1 when all
  electrons are paired, 2 for one unpaired electron (for example low-spin
  iron(III)). Gaussian needs it for every QM region.
\end{description}

\section*{N}

\begin{description}
\item[net charge] The total charge of a molecule, a whole number.
\item[NME] An N-methyl amide cap: a small neutral group added after the last
  residue of a chain.
\item[NPT] A simulation at constant number of particles, pressure and
  temperature. The box size can change, so the water reaches its normal
  density.
\item[N-terminus] The start of a protein chain, with an amino group that
  normally carries a positive charge.
\item[NVT] A simulation at constant number of particles, volume and
  temperature. The box size is fixed.
\end{description}

\section*{O}

\begin{description}
\item[objective] In membrane building, the number packmol lowers while
  packing. It measures how badly molecules still overlap.
\item[occupancy] The share of time an atom spends in a given position, when it
  was seen in more than one.
\item[ONIOM] Gaussian's method for QM/MM calculations, which splits the system
  into layers treated at different levels. See Chapter~\ref{ch:qmmm}.
\item[OPM database] Orientations of Proteins in Membranes: a public collection
  of membrane proteins already placed in a membrane. ProPrep can compare its
  own placement with it.
\item[option] An extra instruction typed after a command, starting with two
  hyphens, such as \texttt{--full-menu}. It applies to that run only. Also
  called a flag. See Chapter~\ref{ch:cmdline}.
\item[orientation] Turning and moving a membrane protein so that its
  water-repelling band lies in the plane of the membrane.
\item[OXT] The second oxygen of the carboxyl group at a charged chain end. It
  exists only on the last residue; many PDB files lack it.
\end{description}

\section*{P}

\begin{description}
\item[packmol-memgen] The program, part of AmberTools, that ProPrep uses to
  pack lipids, water and ions around a membrane protein.
\item[parameters] The numbers in a force field: charges, bond lengths and
  stiffness, angles, twists and atom sizes.
\item[ParmEd] Amber's topology editor. ProPrep uses it to check each new
  topology and to add the 12-6-4 ion terms.
\item[partial charge] A fraction of an electron's charge placed on each atom
  by the force field.
\item[PDB file] A text file (\file{.pdb}) that lists every atom of a structure
  with its position. Most ProPrep steps read and write this format.
\item[PDB format] The older structure file format (\file{.pdb} files), which
  most of ProPrep expects.
\item[PDBQT] The file format AutoDock needs. ProPrep writes it for you from
  the protein and the ligand.
\item[periodic boundaries] The simulation box is repeated in every direction
  like tiles, so a molecule leaving one side enters from the other.
\item[personal access token] A long code, like a password, made on the GitHub
  website, that lets ProPrep's feedback form post on your behalf.
\item[pH] A measure of how acidic the surroundings are. It decides whether
  groups such as histidine or glutamate carry a proton, and so a charge.
\item[pKa] The pH at which half of the molecules of a group carry their
  proton. Below it the group is mostly protonated, above it mostly not.
\item[pLDDT] AlphaFold's confidence score for each residue, from 0 to 100.
  Above 90 is very reliable; below 50 is usually disordered.
\item[plugins] Add-on packages, installed separately, that add their own items
  to ProPrep's menus.
\item[Poisson--Boltzmann (PB)] A way of calculating the electrostatic energy of
  a molecule in salt water. ProPrep's PB titration uses it to predict pKa
  values.
\item[POPC] A common neutral lipid of animal cell membranes, the usual first
  choice for a simple membrane.
\item[port] The number after the colon in a local web address such as
  \texttt{http://127.0.0.1:8000/}. It tells the browser which program on your
  computer to talk to.
\item[positional restraint] A spring-like force that pulls chosen atoms back
  towards their starting positions.
\item[PPM3] A program that works out how a protein sits in a membrane; the
  method behind the OPM database. Slower than MEMEMBED.
\item[principal axes] The natural long, middle and short directions of a
  molecule, worked out from how its mass is spread.
\item[project folder] The folder in which ProPrep saves every file of one
  project. You name it at the start. See Chapter~\ref{ch:projectfolder}.
\item[propionates] The two acid side groups of a heme. Each can be charged or
  neutral. ProPrep's heme sets make them separate residues.
\item[PROPKA] A program that estimates the pKa of each titratable residue from
  the protein structure. See Chapter~\ref{ch:protonation}.
\item[Protein Data Bank (PDB)] The public archive of experimentally solved
  structures of proteins and nucleic acids. Each entry has a four-character
  code such as 1UBQ.
\item[protocol] The ordered list of simulation stages (minimization, heating,
  equilibration, production) and their settings.
\item[protonation state] Whether a titratable group carries its proton, and
  so its charge.
\end{description}

\section*{Q}

\begin{description}
\item[QM (quantum mechanics)] Calculating the electrons explicitly. It can
  describe bond breaking, charge transfer and spin, but is slow and limited
  to a few hundred atoms.
\item[QM/MM] A calculation that treats a small region, such as a redox site,
  with QM and the rest of the protein with MM. See Chapter~\ref{ch:qmmm}.
\item[quaternary structure] How several chains assemble into one molecule,
  for example the four chains of hemoglobin.
\end{description}

\section*{R}

\begin{description}
\item[radius of gyration (Rg)] The average distance of the atoms from the
  center of the molecule; a measure of how compact it is.
\item[receptor] The protein a molecule is docked into.
\item[redocking] Removing a molecule from a structure solved with it, and
  docking it back. It tests whether the settings reproduce a known answer.
\item[redox center] The atom or group in a protein that gains or loses
  electrons, such as the iron of a heme or an iron--sulfur cluster.
\item[redox site] A redox center together with the residues that bind it and
  the bonds between them, as recorded by the Redox Site Detector.
\item[reference] In a superposition, the structure that stays in place while
  the others are moved onto it.
\item[REMARK 350] Lines in a PDB file that describe how to build the
  biological assembly from the asymmetric unit.
\item[REMARK 465] Lines in a PDB file that list residues that were in the
  sample but not seen.
\item[REMARK 470] Lines in a PDB file that list atoms that were not seen.
\item[residue template] The description of one residue in a force field: its
  atoms, their types and charges.
\item[resolution] How fine the detail of an X-ray or cryo-EM structure is, in
  Å. Smaller is better; 2\,Å or less is good.
\item[RESP] Restrained electrostatic potential charges: partial charges
  fitted to a quantum-chemical calculation (in ProPrep, from Gaussian). The
  standard careful method for Amber.
\item[restraints] Extra forces that hold chosen atoms near a position or at a
  distance, for example to protect the protein while the system settles.
\item[RMSD (root-mean-square deviation)] The typical distance, in Å, between
  matching atoms of two structures after lining them up. Near 0 means
  practically identical; in docking, under about 2\,Å from the crystal pose
  counts as a success.
\item[RMSF (root-mean-square fluctuation)] How far each residue wanders around
  its own average position during a simulation. It shows which parts are
  flexible.
\item[route line] The line in a Gaussian input file that starts with
  \texttt{\#} and states the method, basis set and options.
\item[run plan] A saved choice of engines and resources for each stage of a
  protocol, used with a cluster profile.
\end{description}

\section*{S}

\begin{description}
\item[scoring function] The formula that turns a docking pose into a number in
  kcal/mol. ProPrep offers vina, vinardo and ad4. Scores from different
  functions cannot be compared.
\item[scoring matrix] The table BLAST uses to judge how similar two amino
  acids are. BLOSUM62 is the general-purpose default.
\item[search box] The part of the protein that a docking run searches.
\item[SEQRES] Records in a PDB file that give the full sequence of the sample,
  including parts that were not seen.
\item[session file] The file \file{proprep_session_<date>_<time>.json} in
  which ProPrep records every answer you type, in order. It is used to undo,
  resume, replay and make templates. See Chapter~\ref{ch:sessions}.
\item[SLURM] The job system used on most computing clusters: you submit a
  script, and it runs when processors are free.
\item[small model] In the MCPB metal-site method, the small cut-out of the
  metal site, used for the quantum calculation that gives bond and angle
  stiffness.
\item[SMILES] A way of writing a molecule as one line of text, for example
  \texttt{CC(=O)Nc1ccc(O)cc1} for paracetamol.
\item[snapshot] (1)~One saved set of atom positions in a trajectory. (2)~A
  saved copy of your personal force-field library, made with the
  \texttt{proprep-library} command.
\item[SPLIT method] The rule ProPrep uses to count ions: the protein's charge
  is shared between both kinds of ion, so the box is neutral and has the
  chosen salt concentration.
\item[steps] The tiny time increments of a simulation, usually 2~femtoseconds
  each; 500 steps make 1~picosecond.
\item[superposition] Laying one structure on top of another as closely as
  possible, by rotating and moving it.
\end{description}

\section*{T}

\begin{description}
\item[targets] In a superposition, the structures that are moved onto the
  reference.
\item[tautomers] Forms of a molecule that differ only in where a proton sits;
  for histidine, HID and HIE.
\item[template] (1)~The tLEaP input file that ProPrep builds from your
  answers, which you can edit. (2)~In the Redox Site Detector, your own label
  for a detection recipe that is applied to all sites with that label. (3)~A
  recorded session with some answers marked as blanks, so the same steps can
  be run on many proteins (Chapter~\ref{ch:batch}).
\item[TER record] A line in a PDB file that marks the end of a chain or
  molecule. tLEaP needs one after each separate molecule.
\item[Terminal] The program where you type commands to your computer: Terminal
  on macOS, or the Ubuntu window on Windows.
\item[terminal caps] See \emph{cap}.
\item[titratable] Able to gain or lose a proton depending on pH, such as Asp,
  Glu, His, Cys, Tyr, Lys and the chain ends.
\item[tLEaP] The AmberTools program that builds the simulation system: it adds
  missing hydrogens and atoms, applies the force field, adds water and ions,
  and writes the topology and coordinates. See Chapter~\ref{ch:topology}.
\item[TM-align] A program that superimposes proteins by matching residues by
  their position in space, so it works for different sequences with the same
  fold.
\item[TM-score] A score from 0 to 1 for how similar two folds are. Above about
  0.5 means the same fold.
\item[topology] The file (\file{.prmtop}) that lists every atom of the
  simulated system with its charge and type, and every bond, angle and twist
  with its force-field numbers. It holds no positions; those are in the
  coordinate file.
\item[traceback] The full technical error report a program prints when it
  stops with an error: the list of program steps that led to it. See
  Chapter~\ref{ch:troubleshooting}.
\item[trajectory] The file (usually \file{.nc}) of snapshots of all atom
  positions written during a simulation; the main input for analysis.
\item[transformer] A recipe in ProPrep that renames and splits the residues of
  a redox site so that they match a bundled parameter set. See
  Chapter~\ref{ch:redox}.
\item[truncated octahedron] A simulation box with its corners cut off. It fits
  a roughly round protein with less water than a rectangular box.
\end{description}

\section*{U}

\begin{description}
\item[Ubuntu] The Linux system that WSL installs by default. Its window is the
  Terminal you use for ProPrep on Windows.
\item[UniProt] The main database of protein sequences and their descriptions.
\end{description}

\section*{W}

\begin{description}
\item[water layer] In a membrane system, the thickness of water above and
  below the membrane. The default is 17.5\,Å.
\item[water model] The recipe for a single water molecule in a simulation,
  for example TIP3P (three points) or OPC (four points).
\item[Windows Subsystem for Linux (WSL)] A part of Windows 10 and 11 that runs
  a small Linux system alongside Windows. ProPrep runs inside it on Windows.
\item[Workflow Mode] ProPrep's guided main menu, which shows only the tools of
  one stage at a time (Load, Detect, Compare, Fix, Simulate).
\item[workspace] What ProPrep holds in its memory during a run: the
  structures loaded, the files it has made and the results of each tool. It
  starts empty every time ProPrep starts; files on disk stay.
\end{description}
\endgroup
